\documentclass[10pt,a4paper]{amsart}
\usepackage[margin=19mm]{geometry}
\usepackage{amsmath,amssymb,mathtools}
\usepackage[scr=rsfs,cal=euler]{mathalfa}
\usepackage{xfrac}
\usepackage[unicode,hidelinks,bookmarks=false]{hyperref}
\newcommand{\ie}{i.e.\ }
\newcommand{\cf}{cf.\ }
\newcommand{\resp}{respectively\ }
\newcommand{\Subsection}{\subsection}
\providecommand{\nobreakdash}{\nobreak\hbox{-}}
\setlength{\emergencystretch}{2em}
\multlinegap0pt
% ---------------------------------------------------------------------------
% Editorial AMS wrapper prelude. The theorem environments and the mathematical macros below are
% transcribed from the paper's own preamble (cached-originals/source0.tex lines 21--33); the AMS amsart
% class, the named format packages of the pinned TeX Live runtime and this editorial formatting are not
% part of the author's text. No mathematics is changed.
% ---------------------------------------------------------------------------
\newtheorem{theorem}{Theorem}[section]
\newtheorem{definition}[theorem]{Definition}
\newtheorem{lemma}[theorem]{Lemma}
\newtheorem{corollary}[theorem]{Corollary}
\newtheorem{question}[theorem]{Question}
\newtheorem{example}[theorem]{Example}
\newtheorem{proposition}[theorem]{Proposition}
\newtheorem{remark}[theorem]{Remark}
\newtheorem{assumption}[theorem]{Assumption}
\newtheorem{conjecture}[theorem]{Conjecture}



\begin{document}
\title{Characteristic Subgroup Growth}
\author{Liam Hanany and Alexander Lubotzky}
\thanks{The authors were supported by the European Research Council (ERC) under the European Union''s Horizon 2020 (N. 882751).}
\date{Source: arXiv:2510.04150v2; distributed under CC BY 4.0}
\maketitle
\noindent\emph{Source, version and licence.} \emph{Characteristic Subgroup Growth}, by Liam Hanany and Alexander Lubotzky. The source of this editable unit is the e-print of \textbf{version v2} of arXiv:2510.04150, https://arxiv.org/abs/2510.04150v2, whose material is distributed under the Creative Commons Attribution 4.0 International licence, http://creativecommons.org/licenses/by/4.0/, as recorded in the arXiv licence metadata for this paper and shown on that abstract page. The whole of the body below is version v2, the newest version of the paper. Full rights notice: licenses/arxiv-2510.04150-CC-BY-4.0.txt.\par\smallskip
\noindent\textit{Editorial scope.} Counted unit: original Theorem 1.1 of the article, the statement carried by the source environment \texttt{theorem} with source label \texttt{Characteristic Growth of Free Groups}, cached-originals/source0.tex lines 65--68, printed as \textbf{Theorem 1.1} exactly as in the article: for every $r \geq 2$, the characteristic subgroup growth of $F_r$ has type $n^{\mathrm{log}(n)}$. It is delivered together with the author\'s own complete proof as the single primary proof body, source0.tex lines 154--162 (9 lines): the \texttt{proof} environment whose optional argument names it "Proof of Theorem \textbackslash ref\{Characteristic Growth of Free Groups\}", which runs the Zassenhaus-filtration argument, applies the automorphism-action lemma, the exponential-growth corollary and the submodule count, lifts the submodules to characteristic subgroups and concludes the growth type. No proof marker is added and no mathematics is written, repaired or re-derived; the author\'s notation and the printed slips of the source are kept literal. Necessary complete context is retained uncounted, in source order: source0.tex lines 72--75, Theorem 1.2, the companion large-groups statement that the same method proves; lines 90--95, the definition of the Zassenhaus filtration $D_n(\Gamma)$ together with the identity $D_2=\Gamma^p[\Gamma,\Gamma]$; lines 97--100, Lemma 2.2 on the automorphism action (its proof at lines 101--112 is \emph{not} retained because it quotes a book that is not part of this unit); lines 114--119, Lemmas 2.3 and 2.4 on submodule counts; lines 121--124, Corollary 2.5 on $\mathrm{SL}^{\\pm1}_r(\mathbb{F}_p)$-submodules; lines 126--134, Proposition 2.6 computing $\dim(D_n(F_r)/D_{n+1}(F_r))$; and lines 136--151, Corollary 2.7 on exponential growth together with its complete author proof. The article\'s own numbering is reproduced by explicit counter settings: the counter \texttt{theorem}, declared by \texttt{\textbackslash newtheorem\{theorem\}\{Theorem\}[section]} and shared by definition, lemma, corollary, question, example, proposition, remark, assumption and conjecture, is set before each retained passage, so the counted statement prints as Theorem 1.1, the companion statement as Theorem 1.2, the definition as Definition 2.1, the lemmas as Lemmas 2.2, 2.3 and 2.4, the corollaries as Corollaries 2.5 and 2.7 and the dimension computation as Proposition 2.6, exactly as in the article. Every \texttt{\textbackslash ref} and \texttt{\textbackslash eqref} inside every retained passage resolves inside this delivered file. No \texttt{\textbackslash cite} command occurs in any retained passage, so this unit delivers no bibliography entry; the article\'s own bibliography data file \texttt{bibli.bib} is neither spliced into the bundled source nor redistributed. The retained passages contain no figure environment and no graphics-inclusion command, so no artwork is redistributed. The source\'s own class is AMS \texttt{amsart} (\texttt{\textbackslash documentclass[letterpaper]\{amsart\}}); the e-print ships no class or style file, so no publisher class is shipped, used or redistributed, and the wrapper is the same AMS \texttt{amsart} class with the paper\'s own theorem declarations transcribed from its preamble. Every declared body of this unit is an exact, unmodified span of source0.tex: no body is a bounded passage comparison, \texttt{omit} was not used anywhere, and consequently no fidelity receipt is asserted in delivery-evidence.json. The complete concatenated original source is bundled as sources/186420/source0.tex. Omitted from this editable unit and therefore absent from it are source0.tex lines 1--64; 69--71; 76--89; 96--96; 101--113; 120--120; 125--125; 135--135; 152--153; 163--333, that is the article\'s own preamble, title block and abstract, the citation-dense introduction, the proof of Lemma 2.2 with its quoted book, the proofs of Lemmas 2.3 and 2.4 and of Proposition 2.6, the whole section on Segal\'s groups with the counter-example proposition and its proof, the concluding question, and the article\'s bibliography data, all of which remain present in the bundled source but are not delivered here. The AMS \texttt{amsart} wrapper, the named format packages of the pinned TeX Live runtime and this editorial source, licence and scope paragraph are editorial formatting only.\par\smallskip
\setcounter{section}{1}
\setcounter{equation}{0}
\setcounter{theorem}{0}
\begin{theorem}
\label{Characteristic Growth of Free Groups}
For every $r \geq 2$, the characteristic subgroup growth of $F_r$ has type $n^{\mathrm{log}(n)}$.
\end{theorem}

\setcounter{section}{1}
\setcounter{equation}{0}
\setcounter{theorem}{1}
\begin{theorem}
\label{Characteristic Growth of Large Groups}
For every finitely generated $\Gamma$, if $\Gamma$ is large then the characteristic subgroup growth of $\Gamma$ is of type $n^{\mathrm{log}(n)}$.
\end{theorem}

\setcounter{section}{2}
\setcounter{equation}{0}
\setcounter{theorem}{0}
\begin{definition}
Let $\gamma_n(\Gamma)$ be the lower central series of $\Gamma$. Define $D_n(\Gamma) = \prod_{ip^j \geq n} \gamma_i(\Gamma)^{p^j}$.
Equivalently, this sequence can be defined inductively by $D_1(\Gamma) = \Gamma$ and $$D_n(\Gamma) = D_{\lceil\frac{n}{p}\rceil}(\Gamma)^p \prod_{j + k = n}[D_j(\Gamma), D_k(\Gamma)].$$
\end{definition}
In particular $D_2(\Gamma) = \Gamma^p[\Gamma, \Gamma]$.


\setcounter{section}{2}
\setcounter{equation}{0}
\setcounter{theorem}{1}
\begin{lemma}
\label{Automorphism Action}
Let $\varphi:\Gamma \to \Gamma$ be an automorphism of $\Gamma$ acting trivially on the $p$-elementary abelianization $\Gamma/\Gamma^p[\Gamma,\Gamma] = D_1 / D_2$ of $\Gamma$. Then $\varphi$ acts trivially on the subquotients of the Zassenhaus filtration $D_n / D_{n + 1}$.
\end{lemma}

\setcounter{section}{2}
\setcounter{equation}{0}
\setcounter{theorem}{2}
\begin{lemma}
    Let $G$ be a finite group and $p$ a prime number. Then there is a constant $c > 0$ such that for all $\mathbb{F}_p[G]$-modules $M$ of $\mathbb{F}_p$-dimension $d$, there exist submodules $M_1 \leq M_2 \leq M$ such that $M_2 / M_1$ is the direct sum of at least $cd$ isomorphic simple modules.
\end{lemma}
\begin{lemma}
    Let $G$ be a finite group and $M$ a finite dimensional $\mathbb{F}_p[G]$-module. Then there is some $c>0$ such that $M^n$ contains at least $e^{cn^2}$ submodules.
\end{lemma}

\setcounter{section}{2}
\setcounter{equation}{0}
\setcounter{theorem}{4}
\begin{corollary}
\label{Lots of Submodules}
Fix a finite group $G$ and prime $p$. Then there is a constant $c > 0$ such that every $\mathbb{F}_p[G]$-module of $\mathbb{F}_p$-dimension $d$ has at least $e^{cd^2}$-submodules.
\end{corollary}

\setcounter{section}{2}
\setcounter{equation}{0}
\setcounter{theorem}{5}
\begin{proposition}
\label{Zassenhaus dimension computation}
Let $\Gamma = F_r$ be the free group on $r$ generators.
Define
$$w_n(F_r) = \frac{1}{n} \sum_{m | n} \mu(m) r^{\frac{n}{m}}$$
for $\mu$ the Möbius function. Moreover, if $n = p^k m$ for $p \nmid m$ define
$$c_n(F_r) = \sum_{i=0}^kw_{p^im}(F_r)$$
Then $D_n(F_r)/D_{n+1}(F_r)$ is a vector space of dimension $c_n(F_r)$ over $\mathbb{F}_p$.
\end{proposition}

\setcounter{section}{2}
\setcounter{equation}{0}
\setcounter{theorem}{6}
\begin{corollary}
\label{exponential growth}
Let $G = F_r$ be the free group on $r \geq 2$ generators. Then for sufficiently large $n$,
$$\mathrm{dim}_{\mathbb{F}_p}(D_n(F_r) / D_{n+1}(F_r)) \geq \frac{1}{2n} \cdot r^n$$ and
$$[D_n(F_r): D_{n+1}(F_r)] \geq [G: D_{n + 1}(F_r)]^{\frac{1}{15}}$$
\end{corollary}
\begin{proof}
For the first claim it is sufficient to notice that $$c_n(F_r) \geq w_n(F_r) \geq \frac{1}{n} \cdot (r^n - \sum_{k=0}^{\lfloor\frac{n}{2}\rfloor} r^k) \geq \frac{1}{n}(r^n - r^{\lfloor\frac{n}{2}\rfloor + 1})$$
which is greater than $\frac{1}{2n} \cdot r^n$ for sufficiently large $n$.
For the second claim, we start by giving an upper bound on $w_n(F_r)$
$$w_n(F_r) \leq \frac{1}{n} \cdot (r^n + \sum_{k=0}^{\lfloor\frac{n}{2}\rfloor} r^k) \leq \frac{2}{n} \cdot r^n$$
for every $n \geq 1$. We now claim that $c_n(F_r) < \frac{3}{n} \cdot r^n$ for sufficiently large $n.$ Indeed, $$c_n(F_r) \leq w_n(F_r) + \sum_{k=1}^{\lfloor\frac{n}{2}\rfloor} w_k(F_r)\leq \frac{2}{n} \cdot r^n + 2\cdot r^{\lfloor\frac{n}{2}\rfloor + 1} < \frac{3}{n} \cdot r^n$$
for sufficiently large $n,$ say $n \geq n_0$. Denote $M = \sum_{k=1}^{n_0 - 1}c_k(F_r)$.

We can now bound $$\sum_{k = 1}^{n - 1} c_k(F_r) \leq M + 3 \cdot \sum_{k=n_0}^{n - 1} \frac{r^k}{k} \leq M + 3 \cdot (r^{\lfloor \frac{n}{2} \rfloor + 1} + \frac{1}{\frac{n}{2}} \cdot r^n) < \frac{7}{n}  \cdot r^n < 14 \cdot c_n(F_r)$$ for sufficiently large $n.$ This implies that $c_n(F_r) > \frac{1}{15} \cdot \sum_{k=1}^n c_k(F_r)$ as needed.
\end{proof}

\setcounter{section}{3}
\setcounter{equation}{0}
\setcounter{theorem}{0}
\begin{proof}[Proof of Theorem \ref{Characteristic Growth of Free Groups}]
Let $F_r$ be a free group of rank $r$. Let $D_n$ be its Zassenhaus-filtration. This is a descending sequence of characteristic subgroups of $F_r$.

$D_1 / D_2 \simeq (\mathbb{Z}/p)^r$ is an elementary $p$-group. By Lemma \ref{Automorphism Action}, any automorphism acting trivially on $D_1 / D_2$ also acts trivially on each of the subquotients of the Zassenhaus filtration.
The image of $\mathrm{Aut}(F_r)$ in $\mathrm{Aut}(D_1/D_2)$ is a finite group, specifically the group $\mathrm{SL}^{\pm1}_r(\mathbb{F}_p)$ of $r$-by-$r$ matrices in $\mathbb{F}_p$ of determinant $\pm 1$. Therefore, the action of $\mathrm{Aut}(F_r)$ on $D_n/D_{n+1}$ factors through this fixed finite group.

Denote by $N = [G : D_{n + 1}]$. Then $\mathrm{dim}(D_n / D_{n + 1}) \geq \frac{1}{15}\mathrm{log}_p(N)$ by Corollary \ref{exponential growth}.
Hence by Corollary \ref{Lots of Submodules}, we get that there are at least $e^{\frac{c}{225}\mathrm{log}_p(N)^2} = e^{\Omega_{p, r}(\mathrm{log}(n)^2)}$-submodules for this $\mathrm{SL}_r^{\pm 1}(\mathbb{F}_p)$-module. Each such submodule can be lifted and yields a subgroup $D_{n + 1} \leq H \leq D_n$. As this is a submodule for the $\mathrm{Aut}(F_r)$-action, the subgroup $H$ is characteristic in $F_r.$ Moreover, the index of $H$ is bounded by $N$. This proves the claim.
\end{proof}

\end{document}
